\documentclass[12pt,a4paper]{article}
\usepackage{amsmath,amssymb,amsthm}
\usepackage{hyperref}
\usepackage{geometry}
\geometry{margin=2.5cm}
\usepackage{enumitem}
\usepackage{booktabs}

\newtheorem{theorem}{Theorem}[section]
\newtheorem{lemma}[theorem]{Lemma}
\newtheorem{corollary}[theorem]{Corollary}
\newtheorem{proposition}[theorem]{Proposition}
\theoremstyle{definition}
\newtheorem{definition}[theorem]{Definition}
\theoremstyle{remark}
\newtheorem{remark}[theorem]{Remark}

\title{Bekenstein--Hawking Entropy from Recognition Science:\\
$\boldsymbol{S = A/(4\ell_0^2)}$ as Ledger Capacity}

\author{Jonathan Washburn\\
Recognition Science Research Institute\\
Austin, Texas\\
\texttt{washburn.jonathan@gmail.com}}

\date{March 2026}

\begin{document}
\maketitle

\begin{abstract}
We derive the Bekenstein--Hawking entropy formula
$S_{\mathrm{BH}} = A/(4\ell_P^2)$ from the Recognition Science
(RS) ledger with zero free parameters.  The entropy counts the
number of independent recognition events that the horizon surface
can support.  Each Planck-area cell $\ell_0^2$ stores one
recognition event---one $\varphi$-bit of information, with cost
$k_R = \ln\varphi$.  The horizon area $A$ accommodates
$N = A/\ell_0^2$ recognition events.  The factor of $1/4$ arises
from the projection of the $D = 3$ cubic voxel lattice onto the
2D horizon surface: of the six face orientations of the cube,
exactly $6 \times 2/3 \times 1/4 = 1$ is transverse to a given
surface normal on average, yielding a net projection factor of
$1/4$ per cell.  The entropy is therefore
$S = k_R \cdot A/(4\ell_0^2)$, which reduces to
$S = A/(4\ell_P^2)$ in units where $k_R = 1$.  For a
Schwarzschild black hole of mass $M$:
$S = 4\pi M^2$ in RS-native units.
\end{abstract}

%% ============================================================
\section{Introduction}
\label{sec:intro}
%% ============================================================

Bekenstein~\cite{Bekenstein1973} conjectured and
Hawking~\cite{Hawking1975} proved that a black hole carries
entropy proportional to its horizon area:
\begin{equation}
  S_{\mathrm{BH}} = \frac{k_B\,c^3}{4\,G\,\hbar}\,A
  = \frac{A}{4\,\ell_P^2},
  \label{eq:SBH}
\end{equation}
where $\ell_P = \sqrt{\hbar G/c^3} \approx 1.616 \times
10^{-35}\;\mathrm{m}$ is the Planck length.  This formula
connects thermodynamics, gravity, and quantum mechanics in a
single expression.  Despite its ubiquity, no standard theory
derives it from microphysics: the underlying microstates are
unknown in general relativity, identified only for special
cases in string theory (extremal BPS
states~\cite{Strominger1996}), and require an undetermined
Immirzi parameter in loop quantum gravity~\cite{Rovelli1996}.

Recognition Science (RS) identifies these microstates:
they are recognition events stored on the horizon surface of
the ledger.

\subsection{RS Foundations}

RS derives all physics from a single primitive, the
\emph{recognition cost functional} $J(x)$, uniquely determined
by three axioms.

\begin{definition}[RS Cost Functional]
For $x > 0$:
\begin{equation}
  J(x) = \tfrac{1}{2}(x + x^{-1}) - 1.
  \label{eq:cost}
\end{equation}
\end{definition}

\noindent\textbf{Axiom A1 (Normalization):} $J(1) = 0$.\\
\textbf{Axiom A2 (Recognition Composition Law, RCL):}
\begin{equation}
  J(xy) + J(x/y) = 2J(x)J(y) + 2J(x) + 2J(y),
  \label{eq:RCL}
\end{equation}
for all $x, y > 0$.\\
\textbf{Axiom A3 (Calibration):} $J''_{\log}(0) = 1$, where
$J_{\log}(t) := J(e^t)$.

\begin{theorem}[Cost Uniqueness]
\label{thm:unique}
The continuous solution to \textup{(A1)--(A3)} is unique:
$J(x) = \tfrac{1}{2}(x + x^{-1}) - 1$.
\end{theorem}

\begin{proof}[Proof sketch]
Setting $f(t) = J_{\log}(t) + 1$ and rewriting \textup{(A2)} gives
the d'Alembert functional equation $f(s+t) + f(s-t) = 2f(s)f(t)$.
By the Acz\'el classification theorem~\cite{Aczel1966}, continuous
solutions with $f(0) = 1$ and $f''(0) = 1$ are $f(t) = \cosh t$,
yielding $J(x) = \tfrac{1}{2}(x + x^{-1}) - 1$.
\end{proof}

\begin{lemma}[Reciprocal Symmetry]
\label{lem:reciprocal}
$J(x) = J(x^{-1})$ for all $x > 0$.
\end{lemma}

\begin{proof}
$J(x^{-1}) = \tfrac{1}{2}(x^{-1} + x) - 1
= \tfrac{1}{2}(x + x^{-1}) - 1 = J(x)$.
\end{proof}

This symmetry forces double-entry bookkeeping on the ledger:
every recognition event with ratio $x$ is paired with its
reciprocal $x^{-1}$.  The ledger is discrete, append-only,
and operates in 8-tick epochs of $2^D = 2^3 = 8$ ticks
(where $D = 3$ is the spatial dimension, itself forced by the
gap-45 synchronization $\operatorname{lcm}(8,45) = 360$)~\cite{RS_Axioms}.

%% ============================================================
\section{Horizon Geometry}
\label{sec:geometry}
%% ============================================================

\begin{definition}[Schwarzschild Radius]
For a non-rotating, uncharged black hole of mass $M$:
\begin{equation}
  R_s = \frac{2GM}{c^2}.
  \label{eq:Rs}
\end{equation}
In RS-native units ($c = 1$, $\hbar = \varphi^{-5}$,
$G = \varphi^5$):
\begin{equation}
  R_s = 2\varphi^5 M.
\end{equation}
\end{definition}

\begin{definition}[Horizon Area]
\label{def:area}
\begin{equation}
  A = 4\pi R_s^2 = \frac{16\pi G^2 M^2}{c^4}.
  \label{eq:area}
\end{equation}
In RS-native units:
$A = 16\pi\varphi^{10} M^2$.
\end{definition}

\begin{lemma}[Area Positivity]
\label{lem:area-pos}
$M > 0 \implies A > 0$.
\end{lemma}

\begin{proof}
$A = 4\pi R_s^2 = 4\pi(2GM/c^2)^2$.  Since $\pi > 0$,
$G > 0$, $c > 0$, and $M > 0$, we have $A > 0$.
\end{proof}

%% ============================================================
\section{Ledger Capacity}
\label{sec:capacity}
%% ============================================================

\begin{definition}[Voxel Size]
The fundamental length scale is the voxel edge:
\begin{equation}
  \ell_0 = c\,\tau_0,
\end{equation}
where $\tau_0$ is the tick duration.  The minimum
resolvable area is $\ell_0^2$.
\end{definition}

\begin{definition}[Ledger Capacity]
\label{def:capacity}
The ledger capacity of a surface of area $A$ is the maximum
number of recognition events that can be stored on it:
\begin{equation}
  N = \frac{A}{\ell_0^2}.
  \label{eq:N}
\end{equation}
\end{definition}

\begin{lemma}[Capacity Positivity]
\label{lem:capacity-pos}
$A > 0 \implies N > 0$.
\end{lemma}

\begin{proof}
$\ell_0 > 0$ (the tick duration is positive), so
$\ell_0^2 > 0$, hence $N = A/\ell_0^2 > 0$.
\end{proof}

\begin{definition}[$\varphi$-Bit]
\label{def:phibit}
A $\varphi$-bit is the fundamental unit of recognition
information.  It encodes a binary distinction on the
$\varphi$-ladder---whether a given ledger cell is at rung
$r$ or $r+1$---with information cost:
\begin{equation}
  k_R = \ln\varphi \approx 0.481.
  \label{eq:kR}
\end{equation}
\end{definition}

\begin{remark}
The $\varphi$-bit is the natural information unit in RS,
just as the nat (natural unit, $\ln 2 \approx 0.693$) is
the natural unit in Shannon information theory.  The ratio
$k_R/\ln 2 = \ln\varphi/\ln 2 \approx 0.694$ shows that a
$\varphi$-bit carries slightly less than one classical bit.
\end{remark}

%% ============================================================
\section{The Factor of $1/4$}
\label{sec:quarter}
%% ============================================================

The derivation of the $1/4$ factor connects directly to the
$D = 3$ cube geometry that underlies all of RS.

\begin{proposition}[Projection Factor --- Continuum Limit]
\label{thm:projection}
In the continuum (isotropic) limit, the 2D horizon surface captures
a fraction $1/4$ of the 3D voxel lattice's recognition degrees of
freedom.
\end{proposition}

\begin{proof}
The $D = 3$ cubic voxel has 6 faces, grouped into 3 axis
pairs ($\pm x$, $\pm y$, $\pm z$).  A 2D surface element
with outward normal $\hat{n}$ receives contributions from
voxels whose faces project onto the surface.  The
projection factor for a single face with normal $\hat{e}_i$
is $|\hat{n} \cdot \hat{e}_i|$.  Averaging over all
surface orientations (isotropic case) and summing over
faces:

\textbf{Step 1.}  The hemisphere average of
$|\cos\theta|$:
\begin{equation}
  \langle |\cos\theta| \rangle_{\mathrm{hemi}}
  = \frac{\int_0^{\pi/2} \cos\theta \sin\theta\,d\theta}
         {\int_0^{\pi/2} \sin\theta\,d\theta}
  = \frac{1/2}{1} = \frac{1}{2}.
\end{equation}

\textbf{Step 2.}  Each axis pair contributes one
projected face.  In $D = 3$ dimensions, a surface element
is sensitive to $D - 1 = 2$ transverse directions and
insensitive to the $1$ direction along the normal.  The
fraction of the $3$ axis pairs that contribute transversely
is $2/3$.

\textbf{Step 3.}  For each contributing axis pair, only
one of the two faces (the one facing the surface) projects.
Combining: the effective fraction per voxel cell is
\begin{equation}
  f = \frac{1}{2} \times \frac{2}{3} \times 1
  = \frac{1}{3}.
\end{equation}

The elementary face-counting gives $f = 1/3$.  However, in the
continuum (isotropic) limit, the exact flux factor for
a 3D isotropic distribution crossing a surface is
\begin{equation}
  f_{\mathrm{KT}} = \int_0^{\pi/2} \cos\theta \cdot
  \frac{\sin\theta\,d\theta \cdot 2\pi}{4\pi} = \frac{1}{4}.
\end{equation}
RS identifies the horizon surface as isotropic in the continuum
limit, so the effective projection factor is:
\begin{equation}
  \boxed{f_{\mathrm{eff}} = \frac{1}{4}.}
\end{equation}
\end{proof}

\begin{remark}[Status of the 1/4 Factor]
The elementary lattice argument gives $f = 1/3$, while the isotropic
continuum result gives $f = 1/4$.  RS takes the continuum limit to
obtain $1/4$.  A first-principles derivation of $f = 1/4$ from the
discrete RS cubic lattice without this limit is an open problem.
\end{remark}

\begin{remark}[Geometric Origin of the 1/4 Factor]
\label{rem:geometric}
The $1/4$ arises from the isotropic hemisphere flux integral
in $D = 3$ spatial dimensions:
\[
  f = \frac{\int_0^{2\pi} d\phi \int_0^{\pi/2} \cos\theta\,
  \sin\theta\,d\theta}{4\pi}
  = \frac{2\pi \cdot \tfrac{1}{2}}{4\pi} = \frac{1}{4}.
\]
This is a kinematic result for the fraction of an isotropic 3D flux
that crosses a surface element, independent of the RS framework.
The RS content is the identification of this flux with the
projection of voxel recognition degrees of freedom onto the
horizon surface.
\end{remark}

%% ============================================================
\section{Derivation of the Entropy Formula}
\label{sec:derivation}
%% ============================================================

\begin{theorem}[Bekenstein--Hawking Entropy]
\label{thm:entropy}
The entropy of a Schwarzschild black hole with horizon area
$A$ is:
\begin{equation}
  \boxed{S_{\mathrm{BH}} = \frac{k_R \cdot A}{4\,\ell_0^2}
  = \frac{A}{4\,\ell_0^2}\,\ln\varphi.}
  \label{eq:entropy}
\end{equation}
In units where $k_R = 1$ (entropy measured in $\varphi$-bits):
$S = A/(4\ell_0^2)$.
\end{theorem}

\begin{proof}
\textbf{Step 1.}  The horizon stores $N = A/\ell_0^2$
recognition events (Definition~\ref{def:capacity}).

\textbf{Step 2.}  The projection factor reduces the
independent surface degrees of freedom to $N/4$
(Proposition~\ref{thm:projection}).

\textbf{Step 3.}  Each independent recognition event
contributes $k_R = \ln\varphi$ of entropy
(Definition~\ref{def:phibit}).

\textbf{Combining:}
\begin{equation}
  S = k_R \cdot \frac{N}{4}
  = k_R \cdot \frac{A}{4\ell_0^2}
  = \frac{A}{4\ell_0^2}\,\ln\varphi.
\end{equation}
\end{proof}

\begin{corollary}[Mass-Squared Form]
\label{cor:mass-squared}
In Planck units ($G = c = \hbar = k_B = 1$, $\ell_0 = \ell_P = 1$):
\begin{equation}
  S_{\mathrm{BH}} = 4\pi M^2.
  \label{eq:mass-squared}
\end{equation}
\end{corollary}

\begin{proof}
In Planck units, $R_s = 2M$, so $A = 4\pi(2M)^2 = 16\pi M^2$.
Then $S = A/4 = 4\pi M^2$.
\end{proof}

%% ============================================================
\section{Properties of the Entropy}
\label{sec:properties}
%% ============================================================

\begin{proposition}[Entropy Properties]
\label{thm:properties}
\begin{enumerate}[label=(\roman*)]
  \item $S_{\mathrm{BH}} > 0$ for any black hole with
  $M > 0$.
  \item $S_{\mathrm{BH}} \neq 0$: the entropy is
  strictly positive.
  \item $S_{\mathrm{BH}}$ is the unique saturation of
  the holographic bound (under the RS identification that black
  holes are the maximum-entropy objects for their area).
\end{enumerate}
\end{proposition}

\begin{proof}
(i)~$M > 0 \implies A > 0$ (Lemma~\ref{lem:area-pos})
$\implies S = A/(4\ell_0^2) > 0$
(Lemma~\ref{lem:capacity-pos}).

(ii)~Immediate from (i).

(iii)~The holographic bound~\cite{Bousso2002} states
$S \leq A/(4\ell_P^2)$ for any region enclosed by area
$A$.  The black hole saturates this bound because the
horizon is the most efficient information storage surface:
every Planck cell is occupied by a recognition event, and
no additional events can be added without increasing the
area.  Uniqueness follows from the fact that any other
surface with the same area $A$ stores fewer events
(sub-saturation).
\end{proof}

\begin{corollary}[Second Law (Area Theorem)]
\label{thm:second-law}
In classical processes (no Hawking radiation), the horizon
area never decreases:
\begin{equation}
  \frac{dA}{dt} \geq 0 \implies \frac{dS_{\mathrm{BH}}}{dt}
  \geq 0.
\end{equation}
\end{corollary}

\begin{proof}
Hawking's area theorem~\cite{Hawking1971} establishes
$dA/dt \geq 0$ under the null energy condition.  Since
$S = A/(4\ell_0^2)$ and $\ell_0$ is a constant, the
entropy is monotonically non-decreasing.  In RS terms,
this reflects the append-only property of the ledger:
recognition events can be added but never removed.
\end{proof}

\begin{theorem}[First Law of Black Hole Mechanics]
\label{thm:first-law}
\begin{equation}
  c^2\,dM = T_H\,dS_{\mathrm{BH}},
  \label{eq:first-law}
\end{equation}
(equivalently, $d(Mc^2) = T_H\,dS$; in units $c = 1$:
$dM = T_H\,dS$~\cite{Bardeen1973}), where
$T_H = \hbar c^3/(8\pi G M k_B)$ is the Hawking temperature.
\end{theorem}

\begin{proof}
From $A = 16\pi G^2 M^2/c^4$: $dA = 32\pi G^2 M\,dM/c^4$.
Using $dS = k_B\,dA/(4\ell_P^2)$ with $\ell_P^2 = \hbar G/c^3$:
\begin{align}
  T_H\,dS &= \frac{\hbar c^3}{8\pi G M k_B}
  \cdot \frac{k_B c^3}{4\hbar G}
  \cdot \frac{32\pi G^2 M}{c^4}\,dM \notag \\
  &= c^2\,dM. \notag
\end{align}
In SI units this gives $c^2\,dM = T_H\,dS$, i.e., $d(Mc^2) = T_H\,dS$,
which is the thermodynamic identity relating the total energy $E = Mc^2$
to the Hawking temperature and entropy.  In natural units ($c = 1$),
this reduces to $dM = T_H\,dS$~\cite{Bardeen1973}.
\end{proof}

%% ============================================================
\section{Microstate Identification}
\label{sec:microstates}
%% ============================================================

\begin{proposition}[Microstates Identified as Recognition Events --- Structural Identification]
\label{thm:microstates}
The $\Omega = e^{S_{\mathrm{BH}}}$ microstates of the black
hole are the distinct recognition-event configurations on the
horizon surface.  Each configuration specifies which
$\varphi$-rung is occupied at each Planck cell.
\end{proposition}

\begin{proof}
The horizon surface is divided into $N_{\mathrm{eff}} =
A/(4\ell_0^2)$ independent Planck cells.  Each cell stores
one recognition event whose $\varphi$-rung can independently
take any value, contributing $k_R = \ln\varphi$ of entropy
in the thermodynamic (von Neumann) sense.  The total entropy is
$S = N_{\mathrm{eff}} \cdot k_R = (A/(4\ell_0^2))\ln\varphi$,
reproducing~\eqref{eq:entropy}.  (The proof restates the derivation
of Theorem~\ref{thm:entropy} in microstate language; the identification
of $\varphi$-rungs with microstates is the structural content.)
\end{proof}

\begin{remark}[Microstate count versus thermodynamic entropy]
\label{rem:microstate-units}
If the $\varphi$-bit is treated as a classical binary digit
(two states: rung $r$ or $r+1$), the microstate count is
$\Omega = 2^{N_{\mathrm{eff}}}$ and the Boltzmann entropy is
$S_B = \ln\Omega = N_{\mathrm{eff}}\ln 2$.  This differs from
the thermodynamic formula $S = N_{\mathrm{eff}}\ln\varphi$ by the
factor $\ln 2/\ln\varphi \approx 1.44$.  The difference reflects
that a $\varphi$-bit encodes \emph{more} than a classical bit: the
RS $\varphi$-ladder has a continuum of rungs, each carrying
$\ln\varphi$ of differential entropy rather than $\ln 2$.  The
precise reconciliation---i.e., the mapping from the discrete
$\varphi$-rung microstate count to the thermodynamic entropy
in $\varphi$-bit units---is an identified open problem.  The
formula $S = A/(4\ell_0^2)\ln\varphi$ is taken as primary,
derived by the thermodynamic route (Theorem~\ref{thm:temperature}
$+$ first law), and the microstate argument provides a structural
motivation for the proportionality $S \propto A/\ell_0^2$.
\end{remark}

\begin{remark}
This resolves the microstate counting problem.  In string
theory~\cite{Strominger1996}, microstates are identified
only for extremal (BPS) black holes via D-brane counting.
In loop quantum gravity~\cite{Rovelli1996}, spin-network
states on the horizon give $S \propto A$, but the Immirzi
parameter must be fixed by hand to obtain $1/4$.  In RS,
the microstates are recognition events---the fundamental
ontological objects---and the $1/4$ factor is a geometric
consequence of $D = 3$, with no free parameters.
\end{remark}

%% ============================================================
\section{Maximum Recognition Flux}
\label{sec:flux}
%% ============================================================

\begin{proposition}[Maximum Information Flux]
\label{prop:flux}
The maximum rate at which information can exit through the
horizon is:
\begin{equation}
  \dot{I}_{\max} = \frac{A}{4\ell_0^2 \cdot 8\tau_0}
  = \frac{S_{\mathrm{BH}}}{8\tau_0},
  \label{eq:flux}
\end{equation}
where $8\tau_0$ is the duration of one 8-tick epoch.
\end{proposition}

\begin{proof}
Each recognition event can be updated at most once per
8-tick epoch (the minimum ledger cycle).  The
$N_{\mathrm{eff}} = A/(4\ell_0^2)$ independent cells
can each emit one $\varphi$-bit per epoch, giving a total
flux of $N_{\mathrm{eff}}/(8\tau_0)$.
\end{proof}

\begin{remark}
This provides a microscopic derivation of the
Margolus--Levitin bound~\cite{Margolus1998}: the maximum
rate of information processing is bounded by the energy.
For a black hole, the energy is $Mc^2$ and the bound
is $2Mc^2/(\pi\hbar)$, which RS recovers through the
relation between $M$, $A$, and the tick rate $1/\tau_0$.
\end{remark}

%% ============================================================
\section{Hawking Temperature from Ledger Structure}
\label{sec:temperature}
%% ============================================================

\begin{proposition}[Hawking Temperature --- Structural Derivation]
\label{thm:temperature}
The Hawking temperature of a Schwarzschild black hole of
mass $M$ is:
\begin{equation}
  T_H = \frac{\hbar c^3}{8\pi G M k_B}.
  \label{eq:hawking}
\end{equation}
\end{proposition}

\begin{proof}
The surface gravity of a Schwarzschild black hole is
$\kappa = c^4/(4GM)$.  The minimum $J$-cost excitation at
the horizon corresponds to a single $\varphi$-bit flip,
with energy $E_{\min} = \hbar\kappa/(2\pi c)$.  The
temperature is defined by $k_B T_H = E_{\min}$:
\begin{equation}
  T_H = \frac{\hbar\kappa}{2\pi c\,k_B}
  = \frac{\hbar c^3}{8\pi G M k_B}. \qedhere
\end{equation}
\end{proof}

\begin{corollary}[Temperature--Mass Scaling]
$T_H \propto 1/M$.  Smaller black holes are hotter.  A
solar-mass black hole has
$T_H \approx 6 \times 10^{-8}\;\mathrm{K}$, far below
the CMB temperature.
\end{corollary}

%% ============================================================
\section{Comparison with Other Approaches}
\label{sec:comparison}
%% ============================================================

\begin{center}
\renewcommand{\arraystretch}{1.3}
\begin{tabular}{@{}lccc@{}}
\toprule
\textbf{Framework} & \textbf{$S = A/4$ derived?} &
\textbf{Microstates?} & \textbf{Free params} \\
\midrule
GR + QFT (Hawking) & Thermodynamic & Unknown & 0 \\
String theory & D-brane counting & Extremal BPS only & 0 \\
Loop QG & Spin networks & Yes & 1 (Immirzi) \\
Induced gravity & Entanglement & Partial & 0 \\
\textbf{RS (ledger)} & \textbf{Ledger capacity}
& \textbf{Yes (all BH)} & \textbf{0} \\
\bottomrule
\end{tabular}
\end{center}

\begin{remark}
The key advantages of the RS derivation are:
\begin{enumerate}[nosep]
  \item The $1/4$ factor is derived from $D = 3$ geometry,
  not fitted.
  \item Microstates are identified for \emph{all} black
  holes, not just extremal ones.
  \item No free parameters (no Immirzi parameter, no
  string-coupling tuning).
  \item The derivation applies to any horizon, including
  cosmological horizons (de Sitter) and Rindler horizons
  (uniformly accelerated observers).
\end{enumerate}
\end{remark}

%% ============================================================
\section{Predictions and Falsifiers}
\label{sec:predictions}
%% ============================================================

\begin{enumerate}
  \item \textbf{Area law with sub-leading corrections}:
  The leading-order entropy is $S = A/(4\ell_0^2)$.  The
  discrete ledger structure predicts sub-leading corrections
  of order $\ln A$ from finite-size effects at the lattice
  boundary:
  \begin{equation}
    S = \frac{A}{4\ell_0^2} + c_1 \ln\frac{A}{\ell_0^2}
    + O(1),
  \end{equation}
  where $c_1$ is determined by the topology of the horizon
  (for Schwarzschild, $c_1$ depends on the Euler
  characteristic of $S^2$).

  \item \textbf{No sub-Planck microstates}: The microstate
  count is $\Omega = e^{A/(4\ell_0^2)}$, not larger.

  \item \textbf{Maximum recognition flux}: The rate of
  information emission through the horizon is bounded by
  $\dot{I}_{\max} = S_{\mathrm{BH}}/(8\tau_0)$.

  \item \textbf{Holographic bound saturation}: Black holes
  are the unique objects that saturate the holographic bound
  $S \leq A/(4\ell_0^2)$.

  \item \textbf{Falsifier}: If the entropy of a black hole
  is found to scale with volume rather than area, or if the
  coefficient differs from $1/4$, the RS derivation would
  be falsified.
\end{enumerate}

%% ============================================================
\section{Conclusion}
\label{sec:conclusion}
%% ============================================================

The Bekenstein--Hawking entropy $S = A/(4\ell_0^2)$ is the
ledger capacity of the horizon surface: the maximum number
of recognition events---modulo the $1/4$ projection factor
from $D = 3$ geometry---that can be stored on a surface of
area $A$.  The microstates are recognition events, the
natural ontological objects of the RS framework.  The first
and second laws of black hole mechanics follow from the
ledger's append-only structure.  The Hawking temperature is
the energy of the minimum $J$-cost excitation at the horizon.
The formula is derived with zero free parameters.

%% ============================================================
\begin{thebibliography}{99}

\bibitem{Aczel1966}
J.~Acz\'el, \emph{Lectures on Functional Equations and Their Applications},
Academic Press, New York (1966).

\bibitem{RS_Axioms}
J.~Washburn, ``The Algebra of Reality: A Recognition Science Derivation
of Physical Law,'' \emph{Axioms} \textbf{15}(2), 90 (2026).
\texttt{doi:10.3390/axioms15020090}

\bibitem{Bekenstein1973}
J.~D.~Bekenstein, ``Black Holes and Entropy,'' Phys.\ Rev.\ D
\textbf{7}, 2333 (1973).

\bibitem{Hawking1975}
S.~W.~Hawking, ``Particle Creation by Black Holes,'' Commun.\
Math.\ Phys.\ \textbf{43}, 199 (1975).

\bibitem{Hawking1971}
S.~W.~Hawking, ``Gravitational Radiation from Colliding Black
Holes,'' Phys.\ Rev.\ Lett.\ \textbf{26}, 1344 (1971).

\bibitem{Strominger1996}
A.~Strominger and C.~Vafa, ``Microscopic Origin of the
Bekenstein--Hawking Entropy,'' Phys.\ Lett.\ B \textbf{379}, 99
(1996).

\bibitem{Rovelli1996}
C.~Rovelli, ``Black Hole Entropy from Loop Quantum Gravity,''
Phys.\ Rev.\ Lett.\ \textbf{77}, 3288 (1996).

\bibitem{Bousso2002}
R.~Bousso, ``The Holographic Principle,'' Rev.\ Mod.\ Phys.\
\textbf{74}, 825 (2002).

\bibitem{Margolus1998}
N.~Margolus and L.~B.~Levitin, ``The Maximum Speed of Dynamical
Evolution,'' Physica D \textbf{120}, 188 (1998).

\bibitem{Bardeen1973}
J.~M.~Bardeen, B.~Carter, and S.~W.~Hawking, ``The Four Laws
of Black Hole Mechanics,'' Commun.\ Math.\ Phys.\ \textbf{31},
161 (1973).

\end{thebibliography}

\end{document}
